\section*{Research Question}

How do mandatory paternity leave quotas affect the life-cycle gender wage gap and household welfare? Specifically, does forcing fathers to take a minimum share of parental leave reduce the long-run earnings gap between partners, and at what welfare cost?

\section*{Motivation}

Using Danish register data, \citet{kleven2019children} document a 20-percentage-point ``child penalty'' on mothers' earnings after first birth that does not appear for fathers, with mothers taking approximately 80 percent of total parental leave. \citet{kleven2019childpenalties} show the pattern is pervasive across countries, and \citet{patnaik2019reserving} finds that earmarked paternity quotas substantially shift leave toward fathers. Because such quotas are recent --- Denmark's 2022 use-or-lose reform is the leading example --- and their long-run career effects play out over decades, observational data cannot yet identify the causal impact: answering whether the female-earnings gain outweighs the household income loss requires a structural life-cycle model.

\section*{Model}

We propose a finite-horizon two-earner household model in the spirit of \citet{borella2023marriage}, indexed by $i \in \{1,2\}$ (father, mother) over $T$ periods. The household maximises joint expected lifetime utility by choosing consumption, hours, and the leave allocation, subject to a policy constraint on the father's minimum leave share. Table~\ref{tab:params} lists the parameters.

\subsection*{Recursive Formulation}

The state is $(n_t, a_t, k_{1,t}, k_{2,t})$, where $n_t \in \{0,1\}$ indicates the presence of a child, $a_t$ is assets, and $k_{i,t}$ is human capital. The Bellman equation is
\begin{equation}
    V_t(n_t, a_t, k_{1,t}, k_{2,t}) = \max_{c_t,\, h_{1,t},\, h_{2,t},\, \lambda_t}\;\Bigl\{ U(c_t, h_{1,t}, h_{2,t}, n_t) + \rho\,\mathbb{E}_t\!\left[V_{t+1}(\cdot)\right]\Bigr\},
\end{equation}
with period utility
\begin{equation}
    U(c_t, h_{1,t}, h_{2,t}, n_t) = \frac{c_t^{1+\eta}}{1+\eta} - \beta(n_t)\!\left(\frac{h_{1,t}^{1+\gamma}}{1+\gamma} + \frac{h_{2,t}^{1+\gamma}}{1+\gamma}\right),
\end{equation}
and $\beta(n_t) = \beta_0 + \beta_1\, n_t$.

\subsection*{Wages and Human Capital}

Wages are pre-tax and linear in own human capital, with a gender-specific scale that captures the empirical pre-birth wage gap,
\begin{equation}
    w_{i,t} = w_i\,(1 + \alpha\, k_{i,t}),
\end{equation}
and human capital depreciates only with parental leave,
\begin{equation}
    k_{i,t+1} = k_{i,t}\,(1 - \delta_i\, l_{i,t}) + h_{i,t}.
\end{equation}
The asymmetry $\delta_2 \geq \delta_1$ reflects that mothers' leave coincides with the most career-sensitive infant-care phase while fathers' shorter blocks fall in lower-intensity periods, consistent with \citet{kleven2019children}.

\subsection*{Fertility, Leave, and Policy}

A child arrives stochastically with probability
\begin{equation}
    p(n_t) = \begin{cases} p_b & n_t = 0,\\ 0 & n_t = 1.\end{cases}
\end{equation}
Upon birth, total leave $\bar{L}$ is split between partners with father's share $\lambda_t$, each receiving a wage replacement at rate $\phi$. Leave is taken only in the birth period:
\begin{equation}
    l_{1,t} = \begin{cases}\lambda_t \bar{L} & \text{birth in } t\\ 0 & \text{otherwise}\end{cases},\qquad
    l_{2,t} = \begin{cases}(1-\lambda_t) \bar{L} & \text{birth in } t\\ 0 & \text{otherwise}\end{cases}.
\end{equation}
Each agent's time satisfies $h_{i,t} + l_{i,t} \leq 1$, and the paternity quota imposes $\lambda_t \geq \bar{q}$ at birth, where $\bar{q} \in [0, 0.5]$ is the policy instrument. The household budget constraint is
\begin{equation}
    a_{t+1} = (1+r)\!\left(a_t + \sum_{i}\!\bigl[(1-\tau_t)\,w_{i,t}\,h_{i,t} + \phi\, w_i\, l_{i,t}\bigr] - c_t\right).
\end{equation}

\subsection*{Mechanism}

A higher $\bar{q}$ shifts leave from mother to father, lowering her depreciation $\delta_2 l_{2,t}$ and raising her future wage; but because $\delta_1 < \delta_2$, reallocating leave is costly in terms of aggregate human capital and household income. The optimal $\bar{q}$ balances these gains and losses.

\section*{Identification and Counterfactuals}

\paragraph{Calibration.} Preference and technology parameters $\eta, \gamma, \rho, r, \alpha, p_b, \beta_0, \beta_1$ are set to standard values from the literature; the wage scales are normalised; the tax rate $\tau_t$, replacement rate $\phi$, and total leave $\bar{L}$ come from Danish legislation and \emph{Barselstatistik}.

\paragraph{Disciplining the asymmetry.} The two depreciation rates $\delta_1, \delta_2$ generate the child penalty. We discipline them with two targeted moments from \citet{kleven2019children}: $\delta_2$ matches the roughly 20-percentage-point drop in mothers' relative wages after first birth, and $\delta_1$ ensures essentially no father penalty.

\paragraph{Counterfactuals.} With $\theta = (\delta_1, \delta_2)$ pinned down, we re-solve and re-simulate for $\bar{q} \in \{0,\, 0.10,\, 0.25,\, 0.50\}$ and report (i) the life-cycle gender wage gap $G_t = (w_{1,t}-w_{2,t})/w_{1,t}$, (ii) female labor supply, and (iii) the consumption-equivalent welfare gain $\xi(\bar{q})$, defined by
\begin{equation}
    V_0(n_0, a_0, k_{1,0}, k_{2,0};\, \bar{q}) \;=\; \tilde{V}_0(n_0, a_0, k_{1,0}, k_{2,0};\, \bar{q}=0),
\end{equation}
where $\tilde{V}_0$ is the value under $\bar{q}=0$ with consumption scaled by $(1+\xi)$ in every period. Tracing $\xi(\bar{q})$ and $G_t(\bar{q})$ jointly maps the policy frontier between long-run wage convergence and short-run household income loss.

\section*{Data}

\begin{itemize}
    \item \textbf{Child-penalty estimates from \citet{kleven2019children}:} published event-study coefficients on mothers' and fathers' wages provide the calibration targets for $\delta_1, \delta_2$.
    \item \textbf{Danish social-insurance rules:} the statutory replacement rate $\phi \approx 0.9$ (\emph{barselsdagpenge}) and total leave entitlement $\bar{L}$.
    \item \textbf{Barselstatistik (Statistics Denmark):} baseline father-share of leave $\lambda_0 \approx 0.2$.
\end{itemize}
